problem:  
Let \((M^n, g, f)\) be a complete noncompact *expanding gradient Ricci soliton* satisfying  
\[
\operatorname{Ric}_g + \nabla^2 f = \lambda g, \quad \lambda < 0,
\]  
and suppose that \((M,g)\) is *asymptotically conical* (AC) with asymptotic cone \((C(X),dr^2 + r^2 h)\), where the link \((X,h)\) is *Einstein* with \(\operatorname{Ric}_h = (n-2)h\).  

Let \(\mathcal{L}_h\) denote the Lichnerowicz operator acting on symmetric 2-tensors over \((X,h)\):
\[
\mathcal{L}_h(\dot{h}) = \Delta_L \dot{h} + 2(n-2)\dot{h}.
\]
Assume that \(\dot{h}\) is a smooth, trace-free, divergence-free infinitesimal Einstein deformation of \(h\), i.e. \(\mathcal{L}_h(\dot{h}) = 0.\)

**Problem:**  
Does every nontrivial element of \(\ker \mathcal{L}_h\) correspond to a genuine one-parameter family of asymptotically conical expanding gradient Ricci solitons \((M_t, g_t, f_t)\) such that, as \(r\to\infty\),
\[
g_t \sim dr^2 + r^2 (h + t \dot{h})\, ?
\]
Equivalently, are infinitesimal Einstein deformations of the link metric at infinity unobstructed in the nonlinear deformation theory of asymptotically conical expanding solitons, or does the global soliton potential \(f\) introduce analytic obstructions that collapse these directions?

Why is it a "good" problem:  
This problem removes logical inconsistencies and sharpens the analytic question: it asks whether infinitesimal Einstein deformations of the cone cross-section \(X\) correspond to genuine moduli directions for expanding Ricci solitons. It bridges **local geometric analysis** (spectral theory of the Lichnerowicz operator on Einstein manifolds) with **global deformation theory** of nonlinear geometric PDEs (the Ricci soliton equation). Resolving this question would clarify the structure of the moduli space of asymptotically conical expanding solitons—central to understanding singularity models and stability phenomena in Ricci flow. The formulation is simple (“do infinitesimal cone deformations integrate to full solitons?”) but opens a deep analytic–geometric problem connecting weighted elliptic theory, global potential structure, and asymptotic geometry.